\documentclass[11pt,a4paper]{article}
\usepackage{brainteaser}

\begin{document}
\brainteasertitle{15}{Le plus petit morceau}

\begin{statementbox}
On choisit indépendamment et uniformément deux points de coupe sur une corde
de longueur $1$. On obtient ainsi trois morceaux.

Quelle est la longueur moyenne du plus petit morceau?
\end{statementbox}

\begin{hintsbox}
\begin{enumerate}
  \item Noter $L_1,L_2,L_3$ les longueurs des trois morceaux et
  $M=\min(L_1,L_2,L_3)$.
  \item Pour $0\leq t\leq 1/3$, interpréter l'événement $M>t$.
  \item Les triplets $(L_1,L_2,L_3)$ sont uniformément répartis sur le
  triangle $L_1+L_2+L_3=1$, avec $L_i\geq0$.
  \item Comparer l'aire de ce triangle à celle de la région définie par
  $L_1>t$, $L_2>t$ et $L_3>t$.
  \item Utiliser l'identité
  $\E[M]=\int_0^{\infty}\Pp(M>t)\,\mathrm dt$.
\end{enumerate}
\end{hintsbox}

\correction

Notons $X$ et $Y$ les positions des deux coupes. Elles sont indépendantes et
uniformes sur $[0,1]$. Posons $U=\min(X,Y)$ et $V=\max(X,Y)$. Les longueurs
des trois morceaux, dans l'ordre le long de la corde, sont
\[
L_1=U,
\qquad
L_2=V-U,
\qquad
L_3=1-V.
\]
Elles vérifient donc
\[
L_1+L_2+L_3=1,
\qquad L_i\geq0.
\]

Le couple $(U,V)$ a une densité constante égale à $2$ sur le triangle
$0<u<v<1$: chaque paire ordonnée $(u,v)$ peut en effet provenir de
$(X,Y)=(u,v)$ ou de $(X,Y)=(v,u)$. Le changement de variables
\[
(u,v)\longmapsto (L_1,L_2)=(u,v-u)
\]
est linéaire et son jacobien vaut $1$. Les couples $(L_1,L_2)$ sont donc eux
aussi uniformément répartis sur le triangle
\[
\mathcal T
=\{(l_1,l_2):l_1\geq0,\ l_2\geq0,\ l_1+l_2\leq1\},
\]
le troisième morceau étant imposé par $L_3=1-L_1-L_2$.

Posons maintenant
\[
M=\min(L_1,L_2,L_3).
\]

Pour $0\leq t\leq1/3$, l'événement $M>t$ signifie que chacun des trois
morceaux a une longueur supérieure à $t$. Il est impossible lorsque
$t>1/3$, puisque la somme des longueurs vaut $1$.

Pour $0\leq t\leq1/3$, écrivons
\[
L_i=t+L_i',
\]
où $L_i'>0$. Les nouvelles longueurs ont pour somme
\[
L_1'+L_2'+L_3'=1-3t.
\]
\par\noindent
La région admissible est donc un triangle semblable à $\mathcal T$, avec un
rapport de longueurs égal à $1-3t$. Son aire est multipliée par
$(1-3t)^2$. Comme la loi est uniforme, la probabilité cherchée est le rapport
des aires :
\[
\Pp(M>t)=(1-3t)^2.
\]

\paragraph{Rappel: formule des queues.}
Soit $Z$ une variable aléatoire positive. Pour chaque issue $\omega$, la
fonction indicatrice $\mathbf 1_{\{Z(\omega)>t\}}$ vaut $1$ lorsque
$0\leq t<Z(\omega)$, puis $0$. L'aire sous cette fonction est donc exactement
$Z(\omega)$:
\[
Z(\omega)
=\int_0^\infty
\mathbf 1_{\{Z(\omega)>t\}}\,\mathrm dt.
\]
\par\noindent
L'intégrande étant positif, le théorème de Tonelli permet d'échanger
l'espérance et l'intégrale. De plus, l'espérance d'une indicatrice est la
probabilité de l'événement correspondant. Ainsi,
\[
\begin{aligned}
\E[Z]
&=\E\!\left[
  \int_0^\infty \mathbf 1_{\{Z>t\}}\,\mathrm dt
\right]\\
&=\int_0^\infty
  \E\!\left[\mathbf 1_{\{Z>t\}}\right]\mathrm dt\\
&=\int_0^\infty \Pp(Z>t)\,\mathrm dt.
\end{aligned}
\]
On obtient donc la formule des queues
\[
\boxed{
\E[Z]=\int_0^\infty \Pp(Z>t)\,\mathrm dt
=\int_0^\infty \bigl(1-F_Z(t)\bigr)\,\mathrm dt.}
\]
Si $Z\leq a$, l'intégrale peut naturellement être arrêtée en $a$.

Dans notre problème, $0\leq M\leq1/3$. Par conséquent,
\[
\begin{aligned}
\E[M]
&=\int_0^{\infty}\Pp(M>t)\,\mathrm dt\\
&=\int_0^{1/3}(1-3t)^2\,\mathrm dt\\
&=\frac19.
\end{aligned}
\]

On peut aussi vérifier le résultat à partir de la fonction de répartition.
Pour $0\leq t\leq1/3$,
\[
F_M(t)=\Pp(M\leq t)=1-(1-3t)^2,
\]
d'où la densité
\[
f_M(t)=F_M'(t)=6(1-3t).
\]
Un calcul direct donne alors
\[
\E[M]
=\int_0^{1/3}t f_M(t)\,\mathrm dt
=\int_0^{1/3}6t(1-3t)\,\mathrm dt
=\frac19.
\]

\begin{resultbox}
\[
\boxed{\E[M]=\frac19.}
\]
Le plus petit morceau mesure donc en moyenne environ $0{,}1111$, soit
$11{,}11\,\%$ de la corde.
\end{resultbox}

\paragraph{Généralisation.}
Avec $n-1$ points de coupe indépendants et uniformes, les $n$ longueurs sont
uniformément réparties sur le simplexe de dimension $n-1$. Si $M_n$ désigne
la plus petite d'entre elles, la contrainte $L_i>t$ retire une longueur $t$ à
chacun des $n$ morceaux. Le simplexe restant est donc obtenu par une
homothétie de rapport $1-nt$. Son volume de dimension $n-1$ est multiplié par
$(1-nt)^{n-1}$. Ainsi, pour $0\leq t\leq1/n$,
\[
\Pp(M_n>t)=(1-nt)^{n-1}.
\]
Par conséquent,
\[
\boxed{
\E[M_n]
=\int_0^{1/n}(1-nt)^{n-1}\,\mathrm dt
=\frac1{n^2}.}
\]

\medskip
\noindent
Ce résultat suppose que tous les points de coupe sont choisis sur la corde
initiale. Une procédure séquentielle consistant à choisir un morceau puis à le
recouper peut conduire à une autre loi.

\end{document}
